% =====================================================================
%  Netting des sensibilités vega sous le test de variance EBA
%  Note méthodologique de référence — fidèle à l'implémentation
%  (bibliothèque ebanetting + studio Streamlit "Vega Netting Studio")
%
%  Compilation :  pdflatex vega_netting_note.tex   (deux passes)
%                 ou téléverser sur Overleaf (TeX Live).
%  Dépendances : amsmath, amsthm, mathtools, algorithm2e, babel-french,
%                booktabs, enumitem, hyperref  (toutes standard).
% =====================================================================
\documentclass[11pt,a4paper]{article}

\usepackage[utf8]{inputenc}
\usepackage[T1]{fontenc}
\usepackage[french]{babel}
\usepackage{lmodern}
\usepackage{amsmath,amssymb,amsthm,mathtools}
\usepackage{geometry}
\geometry{margin=2.4cm}
\usepackage{float}                       % spécificateur [H] pour les algorithmes
\usepackage[linesnumbered,ruled,vlined,french,onelanguage]{algorithm2e}
\usepackage{booktabs}
\usepackage{enumitem}
\usepackage{array}
\usepackage{xcolor}
\usepackage[hidelinks]{hyperref}

% ----- macros -----
\newcommand{\code}[1]{\texttt{#1}}
\DeclareMathOperator{\Var}{Var}
\DeclareMathOperator{\Cov}{Cov}
\DeclareMathOperator{\tr}{tr}
\DeclareMathOperator{\Vect}{Vect}
\DeclareMathOperator*{\argmin}{arg\,min}
\DeclareMathOperator*{\argmax}{arg\,max}
\newcommand{\R}{\mathbb{R}}
\newcommand{\dsig}{\Delta\sigma}
\newcommand{\dPi}{\Delta\Pi}
\newcommand{\Nt}{\tilde{N}}
\newcommand{\st}{\tilde{s}}
\newcommand{\abs}[1]{\left| #1 \right|}

% ----- environnements de théorèmes -----
\theoremstyle{plain}
\newtheorem{theoreme}{Théorème}[section]
\newtheorem{proposition}[theoreme]{Proposition}
\newtheorem{lemme}[theoreme]{Lemme}
\newtheorem{corollaire}[theoreme]{Corollaire}
\theoremstyle{definition}
\newtheorem{definition}[theoreme]{Définition}
\newtheorem{hypothese}[theoreme]{Hypothèse}
\theoremstyle{remark}
\newtheorem{remarque}[theoreme]{Remarque}
\newtheorem*{impl}{Correspondance avec le code}

% ----- mots-clés des algorithmes (français explicite) -----
\SetKwInput{KwEntree}{Entrées}
\SetKwInput{KwSortie}{Sortie}
\SetKwInput{KwInit}{Initialisation}

\title{\textbf{Netting des sensibilités vega sous le test de variance}\\[4pt]
\large Valorisation prudente EBA --- AVA Market Price Uncertainty\\
Règlement délégué (UE) 2016/101, art.~9(5) \& art.~89\\[6pt]
\normalsize Note méthodologique de référence (fidèle à l'implémentation)}
\author{\textit{Vega Netting Studio} --- bibliothèque \code{ebanetting}}
\date{Juin 2026}

\begin{document}
\maketitle
\begin{abstract}
\noindent
Cette note expose, de la réglementation à l'algorithme, le problème résolu par
l'outil \emph{Vega Netting Studio} : déterminer à quelle granularité de la surface
de volatilité (maturité $\times$ strike) des sensibilités vega signées peuvent être
\emph{nettées} avant d'appliquer l'écart d'incertitude du consensus, pour le calcul
de l'\emph{Additional Valuation Adjustment} (AVA) au titre du \emph{Market Price
Uncertainty} (RTS EBA 2016/101). On formalise le modèle d'incertitude au premier
ordre, le passage des sensibilités granulaires vers les nœuds de test, le netting
comme opérateur d'agrégation, le test de variance et le plancher de conservatisme,
puis les algorithmes : optimisation gloutonne sous budget de variance, architecture
découplée (structure sans portefeuille / niveau par book), stress de corrélation,
diagnostic spectral, et la méthodologie à deux couches du smile. Le fil conducteur
est qu'\textbf{aucun produit de Kronecker / tensoriel n'est jamais assemblé} : toutes
les variances s'évaluent en forme matricielle. Chaque résultat est accompagné de sa
correspondance avec le module qui l'implémente.
\end{abstract}

\tableofcontents
\bigskip\hrule\bigskip

% =====================================================================
\section{Introduction et cadre réglementaire}
% =====================================================================

\subsection{Le problème}
Sous le Règlement délégué (UE) \textbf{2016/101} (RTS EBA sur la valorisation
prudente), l'\emph{approche principale} impose le calcul d'\emph{AVA} par
\emph{exposition de valorisation} à un \textbf{niveau de confiance de 90~\%}
(art.~89). L'\textbf{AVA Market Price Uncertainty} (MPU) capture l'incertitude sur
le prix de sortie d'une position liée à l'incertitude sur ses paramètres de marché.

L'article \textbf{9(5)} et la Q\&A EBA associée autorisent le calcul de cette AVA
sur une \textbf{exposition nettée} plutôt que sur la somme brute des incertitudes
bucket par bucket --- \emph{à condition que l'établissement démontre que les positions
se compensent réellement face à l'incertitude du paramètre}. Pour le vega
(sensibilité à la volatilité implicite), la question opérationnelle est :
\begin{quote}\itshape
à quelle granularité de la surface de volatilité (maturité $\times$ strike) peut-on
sommer les sensibilités signées avant d'appliquer l'écart d'incertitude du consensus,
tout en restant démontrablement conservateur ?
\end{quote}

\subsection{Le principe de réponse}
Le netting n'est pas accordé par défaut. La compensation doit être \emph{démontrée
robuste} sur des agrégations \emph{partielles} et \emph{interprétables}. L'outil
encadre donc tout schéma de netting entre deux bornes (la somme brute et la pleine
diversification, \S\ref{sec:modele}), et n'autorise une agrégation que si elle
passe un \textbf{test de variance} (la perte de fidélité reste sous un budget,
\S\ref{sec:netting}) et un \textbf{plancher de conservatisme} (le niveau d'AVA ne
descend jamais sous la pleine diversification). L'asymétrie réglementaire structure
tout l'édifice :
\begin{quote}\itshape
le test protège contre la \emph{sous}-estimation de l'AVA ; refuser une compensation
valide est permis (conservateur), conserver une compensation invalidée ne l'est pas.
\end{quote}

\subsection{Plan et invariant de calcul}
La note suit le flux de l'outil : modèle d'incertitude (\S\ref{sec:modele}), passage
aux nœuds de test (\S\ref{sec:passage}), netting et test de variance
(\S\ref{sec:netting}), optimisation (\S\ref{sec:optimal}), architecture découplée
(\S\ref{sec:decouple}), diagnostic spectral (\S\ref{sec:spectral}), méthodologie à
deux couches (\S\ref{sec:twolayer}), audit et mise en œuvre (\S\ref{sec:audit}).

\begin{remarque}[Invariant : pas de produit de Kronecker]
La surface est discrétisée en $n = M \times K$ buckets. La covariance des chocs
d'incertitude est une matrice $(n\times n)$, donc $(MK\times MK)$. Sous l'hypothèse
de découplage (\S\ref{sec:decouple-hyp}), \textbf{cette matrice n'est jamais
assemblée} : toute variance se ramène à un produit matriciel ordinaire sur la grille
$(M,K)$ suivi d'une réduction par somme (Propriété~\ref{prop:property1}). C'est un
choix de coût, de mémoire et d'interprétabilité.
\end{remarque}

% =====================================================================
\section{Le modèle d'incertitude au premier ordre}\label{sec:modele}
% =====================================================================

\subsection{P\&L d'incertitude et Propriété 1}
On note $N \in \R^{M\times K}$ la matrice des vegas par bucket,
$N_{mk} = \partial V / \partial \sigma_{mk}$, et $\dsig \in \R^{M\times K}$ le choc
d'incertitude \emph{mid} centré, de matrice d'écarts-types $s \in \R^{M\times K}$
($s_{mk}>0$) et de structure de corrélation $\rho$. Le produit scalaire de Frobenius
est $\langle A, B\rangle = \sum_{m,k} A_{mk} B_{mk}$.

Au premier ordre, l'incertitude de P\&L de valorisation est
\begin{equation}
\dPi = \langle N, \dsig\rangle,
\qquad
\Var(\dPi) = \langle N, \Sigma\, N\rangle,
\label{eq:dpi}
\end{equation}
où $\Sigma = D \rho D$ avec $D = \operatorname{diag}(s)$ (vectorisation implicite).

\begin{proposition}[Forme matricielle de la variance --- Propriété 1]
\label{prop:property1}
Soit $\rho$ donnée sous \emph{forme découplée par axe}
(Hypothèse~\ref{hyp:decouple}) : $\rho_{(m,k),(m',k')} = \rho^{\mathrm{mat}}_{mm'}\,
\rho^{\mathrm{strike}}_{kk'}$. Alors, pour deux expositions $U,V \in \R^{M\times K}$,
en posant $W_V = V \circ s$ (produit de Hadamard),
\begin{equation}
\Cov\big(\langle U,\dsig\rangle,\langle V,\dsig\rangle\big)
= \big\langle\, U\circ s,\ \rho^{\mathrm{mat}}\,(V\circ s)\,\rho^{\mathrm{strike}}\,\big\rangle,
\label{eq:property1}
\end{equation}
un \emph{sandwich} de produits matriciels usuels suivi d'une somme. En particulier
$\Var(\langle V,\dsig\rangle) = \langle V\circ s,\ \rho^{\mathrm{mat}}(V\circ s)
\rho^{\mathrm{strike}}\rangle$. Si une corrélation pleine $\rho^{\mathrm{full}}$
$(n\times n)$ est fournie, on calcule directement
$w_U^\top \rho^{\mathrm{full}} w_V$ avec $w = \mathrm{vec}(V\circ s)$.
\end{proposition}

\begin{impl}
\code{MarketDataBundle.covariance\_of} et \code{variance\_of}
(\code{ebanetting/datasource.py}). La forme \eqref{eq:property1} est
\code{sum(wa * (corr\_mat @ wb @ corr\_strike))} ; le cas plein est
\code{wa.flatten() @ corr\_full @ wb.flatten()}. \code{UncertaintyModel}
(\code{model.py}) précalcule $\Var(\dPi)$.
\end{impl}

\subsection{Les deux extrêmes de l'AVA}
Soit $\kappa$ le multiplicateur de quantile (sous l'hypothèse gaussienne,
$\kappa = z_{0.90} \approx 1{,}2816$, art.~89 ; \code{KAPPA\_90} dans le code). Les
deux régimes extrêmes bornent toute AVA :
\begin{align}
\text{add-up (aucun netting)} &:\quad
\mathrm{AVA}_{\mathrm{brut}} = \kappa \sum_{m,k} \abs{N_{mk}}\, s_{mk},
\label{eq:ava-brut}\\[2pt]
\text{pleine diversification} &:\quad
\mathrm{AVA}_{\mathrm{full}} = \kappa \sqrt{\Var(\dPi)}
\;\le\; \mathrm{AVA}_{\mathrm{brut}}.
\label{eq:ava-full}
\end{align}
L'écart $\eqref{eq:ava-brut}-\eqref{eq:ava-full}$ est le \textbf{bénéfice de netting
théorique maximal}. Le \emph{budget de variance} introduit plus loin est
\begin{equation}
B(\alpha) = (1-\alpha)\,\Var(\dPi), \qquad \alpha \in [0{,}90\,;\,0{,}95].
\label{eq:budget}
\end{equation}

\begin{impl}
\code{UncertaintyModel.ava\_brut}, \code{ava\_full}, \code{budget}
(\code{model.py}). L'inégalité \eqref{eq:ava-full} est l'inégalité de
Cauchy--Schwarz / sous-additivité de l'écart-type.
\end{impl}

% =====================================================================
\section{Hypothèse de découplage et données}\label{sec:decouple-hyp}
% =====================================================================

\begin{hypothese}[Découplage des corrélations par axe]\label{hyp:decouple}
La corrélation 2-D se factorise terme à terme :
$\rho_{(m,k),(m',k')} = \rho^{\mathrm{mat}}_{mm'}\,\rho^{\mathrm{strike}}_{kk'}$, avec
$\rho^{\mathrm{mat}}\in\R^{M\times M}$ (corrélation entre maturités) et
$\rho^{\mathrm{strike}}\in\R^{K\times K}$ (corrélation entre moneyness).
\end{hypothese}

Cette factorisation n'introduit que $\binom{M}{2} + \binom{K}{2}$ paramètres, chacun
interprétable et stressable indépendamment, au lieu de $\binom{MK}{2}$. Elle est
équivalente à une structure de Kronecker $\rho = \rho^{\mathrm{strike}} \otimes
\rho^{\mathrm{mat}}$ sur la vectorisation, mais cette matrice n'est jamais formée :
la Propriété~\ref{prop:property1} l'évite. Une corrélation pleine reste admise
($\rho^{\mathrm{full}}$, prioritaire).

\paragraph{Contrat de données.} L'entrée est un \code{MarketDataBundle} : libellés et
années de maturité, niveaux de strike (en \textbf{moneyness} $K/F$), matrices $N$ et
$s$ de taille $(M,K)$, et corrélations découplées $(\rho^{\mathrm{mat}},
\rho^{\mathrm{strike}})$ ou pleine $\rho^{\mathrm{full}}$. La validation
(\code{validate}) vérifie les dimensions, $s>0$, la symétrie, la diagonale unité et la
semi-définie positivité.

\begin{remarque}[Pourquoi le moneyness, et l'asymétrie sur $s$ et $\rho$]
En strike \emph{absolu}, les points dérivent en moneyness avec le spot : la
corrélation de l'axe strike est gonflée et le netting \textbf{sur-justifié} --- un
biais non conservateur (note \S8). On estime donc $s,\rho$ en coordonnées $x=K/F(T)$.
De même, $s$ est la dispersion inter-contributeurs (Totem) : avec peu de
contributeurs il faut l'\emph{élargir}, pas le lisser ; et $\rho$ s'estime sur les
\emph{variations} des marks de consensus, jamais sur les niveaux (la cointégration
gonfle les corrélations de niveaux). En cas de doute, le stress de corrélation doit
être \emph{adverse au netting}.
\end{remarque}

\begin{impl}
\code{datasource.py} : \code{MarketDataBundle} (contrat, (dé)sérialisation JSON,
\code{validate}), \code{nearest\_correlation} (projection sur le cône des matrices de
corrélation valides), \code{SyntheticDataSource} (books de démonstration),
\code{DataSource}/\code{ScenarioSource} (points de branchement bancaires).
\end{impl}

% =====================================================================
\section{Étape 1 --- le passage aux nœuds de test}\label{sec:passage}
% =====================================================================

L'incertitude du consensus ($s,\rho$) n'est observable qu'à quelques \emph{piliers}
(dispersions Totem, fourchettes brokers). Le vega vit sur une grille plus fine. On le
\emph{transporte} sur la grille de piliers avant tout netting.

\subsection{Matrices de passage et sandwich}

\begin{definition}[Matrice de passage --- Déf. 1]\label{def:passage}
Une matrice de passage 1-D $A \in \R^{g\times q}$ (de la grille fine $g$ vers les $q$
piliers) a des lignes \textbf{positives sommant à un}. Trois conventions remplissent
les lignes pour un point $t$ encadré par les piliers $p_a < t < p_b$ :
\begin{itemize}[nosep]
\item \emph{interp} : poids d'interpolation linéaire $A_{\cdot a}=\frac{p_b-t}{p_b-p_a}$,
      $A_{\cdot b}=\frac{t-p_a}{p_b-p_a}$ ;
\item \emph{equal} : $A_{\cdot a}=A_{\cdot b}=\tfrac12$ ;
\item \emph{quadrant} : tout le poids sur le pilier le plus proche.
\end{itemize}
Hors de la plage des piliers, extrapolation plate (tout le poids sur le pilier extrême).
\end{definition}

\begin{definition}[Passage 2-D (sandwich) --- Déf. 2]
Le vega projeté sur les nœuds de test est
\begin{equation}
\Nt = A_T^\top\, N\, A_K \;\in\; \R^{q_T\times q_K},
\label{eq:sandwich}
\end{equation}
avec $A_T\in\R^{M\times q_T}$ (axe maturité) et $A_K\in\R^{K\times q_K}$ (axe strike) ---
deux produits matriciels ordinaires, un par axe.
\end{definition}

\noindent Le vega est \emph{contracté} par \eqref{eq:sandwich} ; l'incertitude
($s$, $\rho$) est \emph{restreinte} aux indices des piliers (sous-matrices), là où elle
est mesurée. La conservation du vega est immédiate puisque les lignes de $A$ somment à
un :
\begin{equation}
\textstyle\sum_{m,k} N_{mk} \;=\; \sum_{a,b}\Nt_{ab}
\qquad\text{(Propriété 3, conservation du vega).}
\label{eq:conservation}
\end{equation}

\subsection{Le théorème de cohérence passage--surface}

\begin{definition}[Tracking error de passage --- éq. 5]
Soit $B_T,B_K$ les matrices de passage \emph{interp}. Pour une convention de matrices
$A_T,A_K$, l'erreur de représentation est $E = B_T^\top N B_K - A_T^\top N A_K$, et la
tracking error de passage
\begin{equation}
\mathrm{TE}^2_{\mathrm{passage}} = \Var\big(\langle E, \dsig^{\mathrm{pil}}\rangle\big).
\label{eq:te-passage}
\end{equation}
\end{definition}

\begin{theoreme}[Cohérence passage--surface --- Th. 1]\label{th:passage}
Lorsque les poids de passage coïncident avec les \emph{poids d'interpolation de la
construction de la surface de pricing} ($A=B$), le passage ne détruit aucune
information au premier ordre :
\[ \mathrm{TE}^2_{\mathrm{passage}} = 0. \]
Toute autre convention (quadrant, équipondérée) crée une erreur quantifiable
\eqref{eq:te-passage} strictement positive en général, qui \textbf{consomme le budget
de variance \eqref{eq:budget} avant tout netting}.
\end{theoreme}

\begin{proof}[Idée]
Avec $A=B$, $E=0$ par construction, donc \eqref{eq:te-passage} s'annule. Sinon $E\neq0$
et sa variance sous $\Sigma^{\mathrm{pil}}$ est positive sauf annulation accidentelle.
La quantité de référence est $\Var(\langle B_T^\top N B_K, \dsig^{\mathrm{pil}}\rangle)$.
\end{proof}

\begin{remarque}[Les matrices de passage ne sont pas optimisées --- Rem. 3]
$A_T,A_K$ sont \textbf{fixées} par la construction de la surface, ce ne sont pas des
variables d'optimisation. Seule la partition en aval est optimisée. L'outil \emph{mesure}
donc la TE de passage de chaque convention (et avertit si la convention n'est pas
\emph{interp}) au lieu de la minimiser : une TE de passage est une perte sèche de budget,
à documenter et déduire (note \S8).
\end{remarque}

\subsection{Grilles de chocs plus fines : restriction sans perte (\S3.5)}
Quand les chocs sont \emph{plus fins} que les sensibilités (variations quotidiennes de
la surface système, p.\,ex. $350\times100$ contre des piliers $10\times10$), on lit la
surface fine aux piliers par des matrices de sélection $0/1$ :
\begin{equation}
\dsig^{\mathrm{pil}} = R_T\, \dsig^{\mathrm{fin}}\, R_K^\top
\qquad\text{(Propriété 5, exact car $R\,B = I$).}
\end{equation}
La composante hors grille (Remarque 4)
$\dsig^{\mathrm{fin}} - B_T(R_T \dsig^{\mathrm{fin}} R_K^\top) B_K^\top$ est mesurable
quotidiennement : sa taille jauge l'adéquation du jeu de piliers (risque de paramètres
de forme, à traiter à part --- ce n'est pas un défaut du passage). Grilles identiques
$\Rightarrow$ identité.

\begin{algorithm}[H]
\caption{Passage du vega granulaire aux nœuds de test}
\KwEntree{vega $N\in\R^{M\times K}$, indices de piliers $\mathcal{T},\mathcal{K}$,
          convention $c$, modèle ($s,\rho$)}
\KwSortie{bundle de nœuds $(\Nt, s^{\mathrm{pil}}, \rho^{\mathrm{pil}})$, TE de passage
          par convention}
\ForEach{convention $c'\in\{\mathrm{interp},\mathrm{equal},\mathrm{quadrant}\}$}{
  construire $A_T^{(c')}, A_K^{(c')}$ par la Déf.~\ref{def:passage}\;
}
$\Nt \leftarrow (A_T^{(c)})^\top\, N\, A_K^{(c)}$\tcp*{sandwich, éq. \eqref{eq:sandwich}}
restreindre $s,\rho$ aux indices $(\mathcal{T},\mathcal{K})$\;
$N^{\mathrm{ref}} \leftarrow (A_T^{\mathrm{interp}})^\top N\, A_K^{\mathrm{interp}}$\;
\ForEach{convention $c'$}{
  $E \leftarrow N^{\mathrm{ref}} - (A_T^{(c')})^\top N\, A_K^{(c')}$\;
  $\mathrm{TE}^2_{c'} \leftarrow \max\!\big(\Var(\langle E,\dsig^{\mathrm{pil}}\rangle),\,0\big)$
  \tcp*{$=0$ pour interp (Th.~\ref{th:passage})}
}
\Return{nœuds et $\{\mathrm{TE}^2_{c'}\}$}\;
\end{algorithm}

\begin{impl}
\code{passage.py} : \code{passage\_matrix} (Déf.~\ref{def:passage}), \code{project}
(sandwich \eqref{eq:sandwich}), \code{project\_bundle} (algorithme ci-dessus,
\code{te2\_by\_convention}, \code{conserves\_vega}), \code{restriction\_matrix},
\code{restrict\_shocks}, \code{offgrid\_residual} (\S3.5). Onglet \emph{Passage}.
\end{impl}

% =====================================================================
\section{Netting, test de variance et plancher}\label{sec:netting}
% =====================================================================

À partir d'ici on travaille sur la grille de nœuds (notée encore $(M,K)$ pour
alléger, $n = MK$), avec le vega $N$ et le modèle $(\,s,\rho\,)$ restreints aux piliers.

\subsection{Le netting comme opérateur d'agrégation}

\begin{definition}[Schéma de netting --- Déf. 4]\label{def:netting}
Un \emph{schéma} est une partition $\mathcal P = \{\mathcal G_1,\dots,\mathcal G_S\}$
des nœuds (encodée par une matrice de labels entiers $(M,K)$). À chaque ensemble $r$
on associe :
\begin{itemize}[nosep]
\item une \textbf{exposition nettée} $m_r = \sum_{(a,b)\in\mathcal G_r} N_{ab}$ ;
\item une \textbf{grille de poids} $W_r\in\R^{M\times K}$, stochastique en ligne
      ($W_r\ge 0$, support $\subseteq\mathcal G_r$, $\sum W_r = 1$), définissant un
      \textbf{choc représentatif} $\Delta\st_r = \langle W_r, \dsig\rangle$ ;
\item une \textbf{incertitude représentative} $\st_r = \sqrt{\Var(\langle W_r,
      \dsig\rangle)}$ (Propriété~\ref{prop:property1}).
\end{itemize}
Trois pondérations : \emph{pivot} (toute la masse sur la cellule de risque dominant
$\argmax_{(a,b)\in\mathcal G_r}\abs{N_{ab}}\,s_{ab}$), \emph{vega} (proportionnelle à
$\abs{N}$), \emph{equal} (uniforme).
\end{definition}

Le P\&L proxy et le \textbf{résidu} (une matrice $(M,K)$) sont
\begin{equation}
\widehat{\dPi}_{\mathcal P} = \sum_r m_r \langle W_r,\dsig\rangle,
\qquad
R = N - \sum_r m_r W_r .
\label{eq:proxy}
\end{equation}
L'AVA retenue est \textbf{conservatrice} (add-up entre ensembles, valeurs absolues) :
\begin{equation}
\mathrm{AVA}(\mathcal P) = \kappa \sum_r \abs{m_r}\, \st_r .
\label{eq:ava-scheme}
\end{equation}

\subsection{Le test de variance et le plancher}

\begin{definition}[Test de variance et score $R^2$ --- Déf. 3, éq. 6]
La tracking error du schéma est $\mathrm{TE}^2(\mathcal P) = \Var(\langle R,
\dsig\rangle)$. Le test de variance s'écrit
\begin{equation}
\mathrm{TE}^2(\mathcal P) \;\le\; (1-\alpha)\,\Var(\dPi)
\quad\Longleftrightarrow\quad
R^2 := 1 - \frac{\mathrm{TE}^2(\mathcal P)}{\Var(\dPi)} \;\ge\; \alpha .
\label{eq:test}
\end{equation}
\end{definition}

C'est l'analogue prudentiel des tests d'attribution de P\&L (FRTB) : le proxy doit
expliquer une fraction $\alpha$ de la variance du book. Le test \eqref{eq:test}
contrôle la \emph{fidélité} ; le \textbf{plancher de conservatisme} contrôle le
\emph{niveau} :
\begin{equation}
\mathrm{AVA}(\mathcal P) \;\ge\; \kappa\sqrt{\Var(\dPi)} = \mathrm{AVA}_{\mathrm{full}}
\qquad\text{(éq. 7).}
\label{eq:floor}
\end{equation}
Un schéma est \textbf{admissible} ssi il vérifie \eqref{eq:test} et \eqref{eq:floor}.

\begin{algorithm}[H]
\caption{Évaluation d'un schéma de netting}
\KwEntree{labels $(M,K)$, modèle $(N,s,\rho,\kappa)$, seuil $\alpha$, pondération}
\KwSortie{$\mathrm{TE}^2$, $R^2$, $\mathrm{AVA}$, verdicts test/plancher, stats par ensemble}
\For{chaque ensemble $r$}{
  $m_r \leftarrow \sum_{(a,b)\in\mathcal G_r} N_{ab}$\;
  construire $W_r$ selon la pondération (Déf.~\ref{def:netting})\;
  $\st_r \leftarrow \sqrt{\max(\Var(\langle W_r,\dsig\rangle),0)}$\tcp*{Propriété~\ref{prop:property1}}
}
$R \leftarrow N - \sum_r m_r W_r$\;
$\mathrm{TE}^2 \leftarrow \max(\Var(\langle R,\dsig\rangle),0)$\;
$R^2 \leftarrow 1 - \mathrm{TE}^2/\Var(\dPi)$\;
$\mathrm{AVA} \leftarrow \kappa \sum_r \abs{m_r}\,\st_r$\;
$\text{test} \leftarrow [\mathrm{TE}^2 \le (1-\alpha)\Var(\dPi)]$ ;\quad
$\text{plancher} \leftarrow [\mathrm{AVA} \ge \mathrm{AVA}_{\mathrm{full}}]$\;
\Return{$\mathrm{TE}^2, R^2, \mathrm{AVA}, \text{test}, \text{plancher}$}\;
\end{algorithm}

\begin{impl}
\code{netting.py} : \code{NettingScheme} (\code{weight\_grids},
\code{evaluate}$\to$\code{SchemeEvaluation}), \code{SetStat}
(dont le \emph{taux de compensation} $1-\abs{m_r}/\sum\abs{N}$). Onglets
\emph{Netting optimal}, \emph{Labo scénarios}.
\end{impl}

\subsection{La forme fermée à deux nœuds}

Le cas élémentaire éclaire toute la suite : netter le nœud $j$ sur un pivot $i$.

\begin{theoreme}[Cas à deux nœuds --- Th. 2]\label{th:two}
Pour deux nœuds de vegas $\nu_i,\nu_j$, d'incertitudes $s_i,s_j$ et de corrélation
$\rho$, en nettant $j$ sur le pivot $i$ (choc représentatif $s_i$) :
\begin{align}
\textbf{(i)}\quad & \mathrm{TE}^2 = \nu_j^2\,(s_i^2 + s_j^2 - 2\rho\, s_i s_j),
\label{eq:th2-te}\\[2pt]
\textbf{(ii)}\quad & \text{admissible} \iff
\rho \ge \rho_{\min} := \frac{s_i^2 + s_j^2}{2 s_i s_j}
- \frac{(1-\alpha)\,\Var(\dPi)}{2\,\nu_j^2\, s_i s_j},
\label{eq:th2-rhomin}\\[2pt]
\textbf{(iii)}\quad & \text{gain d'AVA} = \kappa\big(\abs{\nu_i}s_i + \abs{\nu_j}s_j
- \abs{\nu_i+\nu_j}\, s_i\big).
\label{eq:th2-gain}
\end{align}
\end{theoreme}

\begin{proof}[Idée]
Le résidu de l'écrasement de $j$ sur $i$ est $R = \nu_j(\delta_j - \delta_i)$ en chocs,
d'où $\Var(\langle R,\dsig\rangle) = \nu_j^2\Var(\dsig_j-\dsig_i) = \nu_j^2(s_i^2+s_j^2
-2\rho s_i s_j)$, soit (i). \eqref{eq:test} donne (ii) en isolant $\rho$.
Enfin l'AVA passe de $\kappa(\abs{\nu_i}s_i+\abs{\nu_j}s_j)$ à
$\kappa\abs{\nu_i+\nu_j}s_i$, d'où (iii).
\end{proof}

\begin{remarque}[Lecture économique]
$\eqref{eq:th2-rhomin}$ dit que netter exige une \emph{corrélation élevée} \textbf{et}
des \emph{incertitudes homogènes} : des points de smile adjacents se nettent, un ATM
1M contre une aile 10Y non. Surtout, plus le vega netté $\abs{\nu_j}$ est grand
relativement au risque total, plus le seuil $\rho_{\min}$ est exigeant : \textbf{on ne
cache pas un gros short derrière un proxy approximatif}. Les fusions de même signe
sont admissibles mais inutiles --- le netting vise des ensembles de signes opposés
fortement corrélés.
\end{remarque}

\begin{impl}
\code{netting.py} : \code{two\_bucket}. Onglet \emph{Labo scénarios} (bac à sable +
frontière $\rho_{\min}(\abs{\nu_j})$).
\end{impl}

% =====================================================================
\section{Netting optimal sous budget de variance}\label{sec:optimal}
% =====================================================================

\subsection{Le problème et sa restriction}
Le problème optimal s'écrit
\begin{equation}
\mathcal P^\star = \argmin_{\mathcal P}\ \kappa \sum_r \abs{m_r}\,\st_r
\quad\text{s.c.}\quad
\mathrm{TE}^2(\mathcal P) \le (1-\alpha)\Var(\dPi)\ \text{et}\ \eqref{eq:floor}.
\label{eq:opt}
\end{equation}
L'espace des partitions a la taille des \emph{nombres de Bell} : \eqref{eq:opt} est
NP-difficile. On le restreint aux \textbf{pavages rectangulaires contigus} de la
grille : chaque ensemble est un rectangle $[r_0,r_1]\times[c_0,c_1]$, et l'on ne
fusionne que deux rectangles \emph{adjacents dont l'union est un rectangle}. Ce choix
est doublement motivé : les rectangles sont interprétables et documentables au titre de
l'art.~9(5) (« maturités 1M--3M $\times$ strikes 0{,}9--1{,}1 »), et il réduit
drastiquement la combinatoire.

\subsection{Le glouton agglomératif (\S7.3)}
On part des singletons ($\mathrm{TE}^2=0$). À chaque étape on évalue toutes les
fusions admissibles (union rectangulaire, $\mathrm{TE}^2$ restant sous le budget) et on
retient la meilleure selon l'\textbf{objectif} choisi :
\begin{itemize}[nosep]
\item objectif \emph{AVA} (défaut) : d'abord les fusions à \emph{coût nul}
      ($\text{coût}\le\varepsilon$, $\text{gain}\ge0$), puis celles maximisant
      $\text{gain}/\text{coût}$, où $\text{gain} = \mathrm{AVA}$ avant $-$ après
      (réduction d'AVA) et $\text{coût} = \mathrm{TE}^2$ après $-$ avant ;
\item objectif \emph{TE seule} (optionnel) : on classe par coût croissant
      ($-\text{coût}$, les moins chères d'abord), \emph{sans} pondérer par l'AVA ---
      utile pour répondre « quelle est la maille la plus grossière justifiable par la
      seule fidélité ? ». La partition obtenue n'est pas garantie optimale en AVA.
\end{itemize}
Enfin (\textbf{étape 3}) on vérifie le plancher \eqref{eq:floor} : tant qu'il est
violé, on défait la dernière fusion (\emph{roll-back}).

\begin{algorithm}[H]
\caption{Glouton agglomératif sous budget (objectifs AVA / TE) + roll-back du plancher}
\KwEntree{modèle $(N,s,\rho,\kappa)$, seuil $\alpha$, pondération, objectif
          $o\in\{\text{ava},\text{te}\}$}
\KwSortie{partition $\mathcal P$, évaluation, historique des fusions, nb de roll-backs}
\KwInit{rectangles $\leftarrow$ singletons ; $\mathit{ev}\leftarrow$ évaluer ;
        $B \leftarrow (1-\alpha)\Var(\dPi)$ ; $\varepsilon \leftarrow 10^{-12}\max(\Var(\dPi),1)$}
\While{$\abs{\text{rectangles}} > 1$}{
  $\text{meilleur} \leftarrow \varnothing$\;
  \For{chaque paire $(i,j)$ de rectangles fusionnables}{
    $\mathit{ev}' \leftarrow$ évaluer la partition d'essai\;
    $\text{gain} \leftarrow \mathit{ev}.\mathrm{AVA} - \mathit{ev}'.\mathrm{AVA}$ ;\quad
    $\text{coût} \leftarrow \mathit{ev}'.\mathrm{TE}^2 - \mathit{ev}.\mathrm{TE}^2$\;
    \If{$\mathit{ev}'.\mathrm{TE}^2 > B$}{\textbf{continuer}\tcp*{hors budget}}
    \uIf{$o = \text{te}$}{$\text{cand} \leftarrow (\text{niveau }0,\ -\text{coût})$}
    \uElseIf{$\text{coût}\le\varepsilon$ \textbf{et} $\text{gain}\ge0$}{
      $\text{cand} \leftarrow (\text{niveau }1,\ \text{gain})$\tcp*{coût nul prioritaire}}
    \uElseIf{$\text{gain}>0$ \textbf{et} $\text{coût}>\varepsilon$}{
      $\text{cand} \leftarrow (\text{niveau }0,\ \text{gain}/\text{coût})$}
    \Else{\textbf{continuer}}
    \lIf{$\text{cand} > \text{meilleur}$}{$\text{meilleur}\leftarrow\text{cand}$}
  }
  \lIf{$\text{meilleur} = \varnothing$}{\textbf{arrêter}}
  appliquer la fusion ; $\mathit{ev}\leftarrow\mathit{ev}'$ ; journaliser\;
}
\While{plancher \eqref{eq:floor} violé \textbf{et} historique non vide}{
  retirer la dernière fusion ; rejouer ; incrémenter le compteur de roll-back\;
}
\Return{$\mathcal P, \mathit{ev}, \text{historique}, \text{roll-backs}$}\;
\end{algorithm}

\begin{remarque}[Pourquoi cet ordre de priorité]
Les fusions \emph{gratuites} (compensation parfaite, variance résiduelle nulle) sont
prises d'abord. Ensuite, l'objectif AVA achète la réduction d'AVA la plus efficace par
unité de budget de variance dépensé --- un ratio coût/bénéfice prudentiel. L'objectif
TE seule retient au contraire des fusions presque gratuites en fidélité même sans gain
d'AVA (vegas de même signe), maximisant l'agrégation sous le budget.
\end{remarque}

\subsection{Stress de corrélation et netting robuste (\S7.3, étape 4)}
Une fusion n'est défendable en revue de modèle que si elle survit à un stress
\emph{adverse au netting}. Le stress abaisse les corrélations hors diagonale,
\begin{equation}
\rho \;\longmapsto\; \max(\rho - \delta,\,-1),\qquad \delta \sim 0{,}1\text{--}0{,}2,
\label{eq:stress}
\end{equation}
suivi d'une reprojection sur le cône des corrélations valides. Le \textbf{netting
robuste} rejoue le glouton sous le modèle de base et sous chaque $\Sigma$ stressé, puis
prend le \textbf{raffinement commun} : deux buckets ne restent ensemble que s'ils le
sont dans \emph{tous} les régimes. Le schéma retenu est réévalué sous le modèle de base.

\begin{algorithm}[H]
\caption{Netting robuste --- raffinement commun sous stress de corrélation}
\KwEntree{bundle, $\alpha$, $\kappa$, amplitudes $\{\delta_1,\dots,\delta_p\}$, objectif}
\KwSortie{schéma robuste, évaluation, runs par régime}
$\text{runs}[0] \leftarrow$ glouton(modèle de base)\;
\For{chaque $\delta$}{$\text{runs}[\delta] \leftarrow$ glouton(modèle stressé par \eqref{eq:stress})}
empiler les labels de tous les runs\;
deux buckets ensemble $\iff$ ensemble dans \emph{tous} les runs\tcp*{raffinement commun}
$\mathcal P_{\text{rob}} \leftarrow$ partition résultante\;
\Return{$\mathcal P_{\text{rob}}$, évaluer$(\mathcal P_{\text{rob}},$ modèle de base$)$, runs}\;
\end{algorithm}

\begin{impl}
\code{optimizer.py} : \code{greedy\_netting}/\code{\_greedy} (objectif
\code{"ava"}/\code{"te"}), \code{\_replay} (roll-back), \code{stress\_bundle}
\eqref{eq:stress}, \code{robust\_netting} (raffinement commun). Pavages :
\code{Rect}, \code{\_mergeable}, \code{\_union}. Onglet \emph{Netting optimal}
(toggles objectif et mode robuste).
\end{impl}

% =====================================================================
\section{Architecture découplée : structure $\perp$ niveau}\label{sec:decouple}
% =====================================================================

Le glouton de \S\ref{sec:optimal} mêle un \emph{objectif} (l'AVA, qui dépend du book,
de $s$ et de $\kappa$) à une \emph{contrainte} (le coût en variance, qui ne dépend que
de $\Sigma$). La \S7.5 les sépare en deux runs : la structure est \emph{lente et
justifiée}, le niveau \emph{rapide et par book}. Une partition qui \emph{bouge avec le
book} est en soi un signal d'alerte en validation.

\subsection{Run 1 --- le dendrogramme sans portefeuille}

\begin{definition}[Distance de risque de base --- Th. 2 (i)]
La distance entre deux nœuds est l'écart-type de leur désaccord de choc :
\begin{equation}
d_{ij} = \sqrt{\Var(\dsig_i - \dsig_j)} = \sqrt{s_i^2 + s_j^2 - 2\rho_{ij}\,s_i s_j},
\label{eq:base-risk}
\end{equation}
avec $\rho_{ij} = \rho^{\mathrm{mat}}_{aa'}\rho^{\mathrm{strike}}_{bb'}$ (découplage) ou
$\rho^{\mathrm{full}}_{ij}$. C'est exactement la racine du coût du Théorème~\ref{th:two}
(i) : fusionner $j$ sur le pivot $i$ coûte $\mathrm{TE}^2 = \nu_j^2\, d_{ij}^2$.
\end{definition}

On agglomère les nœuds (rectangles contigus) par \emph{risque de base croissant}. Le
critère de fusion est le $\varepsilon$ de la Propriété~\ref{prop:prop8} : pour un
ensemble de cellules, le meilleur pivot minimise $\max_j d_{j,\text{pivot}}/s_j$, et
$\varepsilon$ de l'ensemble est ce minimum. Le dendrogramme \textbf{ne dépend que de
$\Sigma$, jamais du book}.

\begin{proposition}[Borne portfolio-free --- Prop. 8]\label{prop:prop8}
Si chaque nœud $j$ d'un ensemble vérifie $d_{j,\text{pivot}} \le \varepsilon\, s_j$,
alors pour \textbf{n'importe quel} book
\begin{equation}
\mathrm{TE} \;\le\; \varepsilon \sum_j \abs{\nu_j}\, s_j
\;=\; \varepsilon\, \frac{\mathrm{AVA}_{\mathrm{brut}}}{\kappa}.
\label{eq:prop8}
\end{equation}
\end{proposition}

\begin{proof}[Idée]
Le résidu d'un ensemble écrasé sur son pivot est $\sum_j \nu_j(\dsig_j -
\dsig_{\text{pivot}})$. L'inégalité triangulaire sur l'écart-type donne
$\mathrm{TE} \le \sum_j \abs{\nu_j}\, d_{j,\text{pivot}} \le \varepsilon\sum_j
\abs{\nu_j} s_j$, puis on somme sur les ensembles.
\end{proof}

\begin{algorithm}[H]
\caption{Run 1 --- construction du dendrogramme (sans portefeuille)}
\KwEntree{bundle $(s,\rho)$}
\KwSortie{arbre de fusions $\{(\text{rect}_a,\text{rect}_b,\text{hauteur }\varepsilon)\}$}
calculer la matrice de distances $d$ \eqref{eq:base-risk}\;
rectangles $\leftarrow$ singletons\;
\While{$\abs{\text{rectangles}} > 1$}{
  \For{chaque paire fusionnable $(i,j)$}{
    $U \leftarrow$ union ; $\varepsilon_U \leftarrow \min_{p\in U}\max_{c\in U}
       d_{c,p}/s_c$\tcp*{critère Prop.~\ref{prop:prop8}}
  }
  fusionner la paire de plus petit $\varepsilon_U$ ; journaliser (hauteur $=\varepsilon_U$)\;
}
\Return{arbre}\;
\end{algorithm}

\subsection{Run 2 --- le niveau par book et le contrôle obligatoire}
Une \emph{coupe} de l'arbre à la hauteur $\varepsilon$ donne une partition. On
l'évalue avec le book du jour, le choc représentatif de chaque ensemble étant son
\textbf{pivot de structure} (issu du Run 1, indépendant du book), $\st_r = s$ au pivot.
Le contrôle obligatoire est le test \eqref{eq:test} ; s'il échoue (ou le plancher), on
\emph{abaisse la coupe} (on défait des fusions) pour ce book. On reporte la borne
Prop.~8 réalisée $\varepsilon_{\text{réel}}\sum\abs{\nu}s$.

\begin{algorithm}[H]
\caption{Run 2 --- coupe et évaluation par book, avec abaissement}
\KwEntree{bundle, $\alpha$, $\kappa$, hauteur $\varepsilon$, dendrogramme}
\KwSortie{labels, évaluation, $\varepsilon_{\text{réel}}$, borne Prop.~8, nb d'abaissements}
$n \leftarrow$ nombre de fusions de hauteur $\le \varepsilon$\;
\Repeat{$(\text{test et plancher})$ \textbf{ou} $n=0$}{
  labels $\leftarrow$ appliquer les $n$ premières fusions\;
  pivots $\leftarrow$ pivots de structure des ensembles\;
  $\mathit{ev} \leftarrow$ évaluer (résidu $=N-\sum_r m_r\,\mathbf 1_{\text{pivot}_r}$)\;
  \lIf{$\neg(\text{test} \wedge \text{plancher})$}{$n \leftarrow n-1$ ; abaissements${+}{+}$}
}
$\varepsilon_{\text{réel}} \leftarrow \max_r \varepsilon_r$ ;\quad
borne $\leftarrow \varepsilon_{\text{réel}}\sum\abs{\nu}s$\;
\Return{labels, $\mathit{ev}$, $\varepsilon_{\text{réel}}$, borne, abaissements}\;
\end{algorithm}

\subsection{Stabilité entre familles de chocs (\S7.4)}
Un même formalisme, plusieurs covariances : $\Sigma^{\mathrm{totem}}$ (niveau d'AVA +
test officiel), $\Sigma^{\mathrm{bid\text{-}ask}}$, $\Sigma^{\mathrm{daily}}$
(variations quotidiennes restreintes aux piliers via la Prop.~5 --- l'estimateur de
$\rho$ le plus dense, \textbf{pas} un niveau prudentiel). On rejoue le test de variance
de la \emph{même} partition sous chaque covariance ; toute fusion qui échoue dans un
régime est \emph{instable} et défaite. Le plancher n'a de sens que sous la covariance
de base.

\begin{algorithm}[H]
\caption{Coupe stable --- survie sous toutes les covariances}
\KwEntree{bundle de base, alternatives $\{\Sigma^{(g)}\}$, $\alpha$, $\kappa$, $\varepsilon$, dendrogramme}
\KwSortie{coupe stable, rapport par régime, nb de fusions défaites}
$n \leftarrow$ nombre de fusions de hauteur $\le \varepsilon$\;
\Repeat{$(\text{ok})$ \textbf{ou} $n=0$}{
  labels $\leftarrow$ coupe à $n$ fusions (pivots de la structure de base)\;
  rapport $\leftarrow$ test de variance sous \emph{chaque} régime\;
  ok $\leftarrow$ (test passé dans \emph{tous} les régimes) \textbf{et} (plancher sous la base)\;
  \lIf{$\neg$ ok}{$n \leftarrow n-1$ ; défaites${+}{+}$}
}
\Return{labels, rapport, défaites}\;
\end{algorithm}

\begin{remarque}[Asymétrie réglementaire, à nouveau]
Refuser une fusion valide est permis ; conserver une fusion invalidée par un régime ne
l'est pas. L'abaissement de coupe (Run 2) et la coupe stable (\S7.4) sont deux
incarnations de ce principe.
\end{remarque}

\begin{impl}
\code{clustering.py} : \code{base\_risk\_distance} \eqref{eq:base-risk},
\code{build\_dendrogram} (Run 1), \code{cut\_pivots}/\code{evaluate\_cut},
\code{decoupled\_netting} (Run 2, \code{te\_bound}, \code{lowered}),
\code{stability\_report}/\code{stable\_cut} (\S7.4). Onglet
\emph{Structure \& Stabilité}.
\end{impl}

% =====================================================================
\section{Diagnostic spectral et décomposition du smile}\label{sec:spectral}
% =====================================================================

\subsection{Facteurs de la covariance (\S6.1)}
La covariance aux nœuds $\Sigma = D\rho D$ (assemblée terme à terme, petite matrice
$q\times q$) se diagonalise dans une base orthonormée (théorème spectral) :
$\Sigma = \sum_\ell \lambda_\ell\, u_\ell u_\ell^\top$, $\lambda_1\ge\cdots\ge\lambda_q\ge0$.

\begin{proposition}[Facteurs et carte de risque --- Prop. 6 et 7]\label{prop:spectral}
Les \emph{facteurs} $\xi_\ell = u_\ell^\top \dsig$ sont décorrélés, de variances
$\lambda_\ell$, et $\sum_\ell \lambda_\ell = \tr(\Sigma) = \sum_i s_i^2$. En posant
$c_\ell = \langle \nu, u_\ell\rangle$ (projection du vega),
\begin{equation}
\dPi = \sum_\ell c_\ell\, \xi_\ell,
\qquad
\Var(\dPi) = \sum_\ell c_\ell^2\, \lambda_\ell .
\label{eq:spectral-var}
\end{equation}
La suite $\{c_\ell^2\lambda_\ell\}$ est la \textbf{carte de risque} du book : elle dit
\emph{où} est le risque, mode par mode.
\end{proposition}

\begin{theoreme}[Plancher spectral --- Th. 3]\label{th:spectral}
Pour tout schéma dont les chocs représentatifs vivent dans
$\Vect(u_1,\dots,u_L)$ --- la description réaliste des schémas grossiers à poids
lisses ---
\begin{equation}
\mathrm{TE}^2 \;\ge\; \sum_{\ell>L} c_\ell^2\,\lambda_\ell .
\label{eq:spectral-floor}
\end{equation}
La variance de \emph{queue} du book est incompressible. On définit le \textbf{nombre
minimal de facteurs} (ou d'ensembles) compatible avec le test :
\begin{equation}
K^\star(\alpha) = \min\Big\{ L :\ \textstyle\sum_{\ell>L} c_\ell^2\lambda_\ell
\le (1-\alpha)\Var(\dPi) \Big\}.
\label{eq:kstar}
\end{equation}
\end{theoreme}

\begin{remarque}[L'hypothèse est nécessaire ; la limite]
Sans l'hypothèse de sous-espace, aucun plancher n'existe : avec $w=\nu$ un seul
ensemble réplique $\dPi$ exactement (Rem. 5). Surtout, \textbf{un facteur n'est pas un
ensemble de netting} : ses poids signés denses ne sont pas une exposition de
valorisation documentable. \emph{Le spectre mesure, la partition nette.}
\end{remarque}

L'inertie $\tau_L = \sum_{\ell\le L}\lambda_\ell / \sum_\ell\lambda_\ell$ (compressibilité
de l'incertitude, U1) ; deux usages complémentaires : \textbf{U3} le \emph{nettoyage
spectral} d'un $\hat\Sigma$ bruité (aplatir la queue des valeurs propres à sa moyenne
en préservant la trace, donc $\sum s_i^2$ ; avec garde-fou de conservatisme) ;
\textbf{U4} la \emph{stabilité des sous-espaces dominants} entre fenêtres (cosinus des
angles principaux $\to 1$).

\begin{algorithm}[H]
\caption{Diagnostic spectral et $K^\star(\alpha)$}
\KwEntree{modèle $(N,s,\rho,\kappa)$, seuil $\alpha$}
\KwSortie{$\{\lambda_\ell\}$, carte de risque $\{c_\ell^2\lambda_\ell\}$, courbe de
          plancher, inertie, $K^\star$}
$\Sigma \leftarrow D\rho D$ ; diagonaliser $(\lambda_\ell, u_\ell)$ par ordre décroissant\;
$c_\ell \leftarrow \langle \nu, u_\ell\rangle$ ;\quad $g_\ell \leftarrow c_\ell^2\lambda_\ell$\;
\For{$L = 0,\dots,q$}{$\text{plancher}[L] \leftarrow \sum_{\ell>L} g_\ell$ ;\quad
   $\tau_L \leftarrow \sum_{\ell\le L}\lambda_\ell/\sum\lambda_\ell$}
$K^\star \leftarrow \min\{L : \text{plancher}[L] \le (1-\alpha)\Var(\dPi)\}$\;
\Return{$\lambda, g, \text{plancher}, \tau, K^\star$}\;
\end{algorithm}

\subsection{Décomposition du smile par tranche (\S6.2)}
Par maturité (tranche $a$), dans la base $\{1, x, x^2\}$ avec $x = K/F - 1$, on lit les
\emph{coordonnées du book} dans la base standard des stratégies de smile :
\begin{equation}
\underbrace{m_a = \textstyle\sum_k \nu_k}_{\text{niveau}},\quad
\underbrace{\mathrm{RR}_a = \textstyle\sum_k \nu_k x_k}_{\text{risk-reversal}},\quad
\underbrace{\mathrm{FLY}_a = \textstyle\sum_k \nu_k x_k^2}_{\text{butterfly}}.
\end{equation}

\begin{theoreme}[Résidu d'écrasement du smile --- Th. 4]\label{th:smile}
Le niveau net $m_a$ se nette \emph{parfaitement quelle que soit sa taille} ; tout le
résidu de l'écrasement de la tranche sur son pivot ATM est porté par $\mathrm{RR}_a$ et
$\mathrm{FLY}_a$. Le go/no-go par tranche est en forme fermée : avec $\Sigma_a = D_a
\rho^{\mathrm{strike}} D_a$ et $e = \nu - m_a\,\mathbf 1_{\text{ATM}}$,
\[ \mathrm{TE}^2_a = e^\top \Sigma_a\, e, \qquad
   R^2_a = 1 - \mathrm{TE}^2_a / (\nu^\top \Sigma_a \nu). \]
\end{theoreme}

Les tranches dont $R^2_a < \alpha$ portent un RR/FLY net matériel : leur forme de smile
doit rester en \emph{add-up} (étape 1 du netting hiérarchique) ; les niveaux des
tranches qui passent poursuivent vers le netting de structure par terme sous
$\rho^{\mathrm{mat}}$.

\begin{impl}
\code{spectral.py} : \code{node\_covariance}, \code{spectral\_diagnostic}
(\eqref{eq:spectral-var}--\eqref{eq:kstar}), \code{spectral\_clean} (U3),
\code{subspace\_stability} (U4). \code{scenario.py} : \code{smile\_decomposition}
(Th.~\ref{th:smile}). Onglet \emph{Spectral \& Smile}.
\end{impl}

% =====================================================================
\section{Méthodologie à deux couches du smile}\label{sec:twolayer}
% =====================================================================

Principe organisateur : le \emph{désaccord} entre points de volatilité est une
propriété du \textbf{sous-jacent} (il ignore le book), tandis que les agrégats RR/FLY
sont la projection du \textbf{book} sur cette structure. D'où deux couches.

\subsection{Couche 1 --- le sous-jacent (sans book)}

\begin{definition}[Modèle de déformation --- Déf. 1]\label{def:smile}
Le choc d'une tranche se décompose en trois modes communs plus un bruit propre :
\begin{equation}
\dsig_j = \dsig_{\mathrm{niv}} + x_j\,\Delta S + x_j^2\,\Delta C + \varepsilon_j,
\end{equation}
d'écarts-types $(s_0, \sigma_S, \sigma_C)$ pour le niveau / la pente (skew) / la
courbure, et $\sigma_\varepsilon$ pour l'idiosyncratique indépendant. L'incertitude
ponctuelle est
$s_x = \sqrt{s_0^2 + x^2\sigma_S^2 + x^4\sigma_C^2 + \sigma_\varepsilon^2}$.
\end{definition}

Le modèle s'estime par régression transversale quotidienne des chocs par strike sur
$[1, x, x^2]$ (les coefficients donnent $(s_0,\sigma_S,\sigma_C)$, le résidu
$\sigma_\varepsilon$, et le $R^2$ poolé la validité).

\begin{theoreme}[Risque de base d'une paire --- Th. 1 (couche 1)]\label{th:tl1}
Le désaccord généré entre deux points est
\begin{equation}
d_{ij}^2 = (x_j-x_i)^2\,\sigma_S^2 + (x_j-x_i)^2(x_j+x_i)^2\,\sigma_C^2
+ 2\sigma_\varepsilon^2 .
\label{eq:tl-dist}
\end{equation}
C'est l'\emph{écart} $\abs{x_j-x_i}$ qui compte (pas la distance à l'ATM) ; le facteur
$(x_i+x_j)^2$ amplifie la courbure dans les ailes ; l'idio est un plancher
$d^2 \ge 2\sigma_\varepsilon^2$.
\end{theoreme}

\begin{proposition}[Corrélation générée --- Prop. 2]
La corrélation de chaque point avec l'ATM n'est pas primitive : le modèle la fabrique,
$\rho(x) = s_0 / s_x$.
\end{proposition}

On construit un \emph{dendrogramme de tranche} : groupes de strikes contigus par
désaccord pivot-à-pivot croissant \eqref{eq:tl-dist}, le pivot d'un groupe étant son
point le plus liquide (plus petit $s$). C'est la \textbf{carte des fusions possibles},
indépendante du book.

\subsection{Couche 2 --- le book (résidu exact)}
Pour un groupe de pivot $p$, on définit le niveau net et les agrégats signés relatifs au
pivot :
\[ m = \textstyle\sum_j \nu_j,\quad
   \mathrm{RR}_p = \textstyle\sum_j \nu_j (x_j - x_p),\quad
   \mathrm{FLY}_p = \textstyle\sum_j \nu_j (x_j^2 - x_p^2). \]

\begin{theoreme}[Résidu de groupe exact --- Th. 2 (couche 2)]\label{th:tl2}
La variance du résidu de l'écrasement du groupe sur son pivot est, \emph{sans
inégalité},
\begin{equation}
\Var(R) = \mathrm{RR}_p^2\,\sigma_S^2 + \mathrm{FLY}_p^2\,\sigma_C^2
+ \sigma_\varepsilon^2\Big(\textstyle\sum_{j\neq p}\nu_j^2
+ \big(\sum_{j\neq p}\nu_j\big)^2\Big).
\label{eq:tl-exact}
\end{equation}
Les modes communs entrent par des sommes \emph{signées} (compensation possible) ;
l'idio par des \emph{carrés} (jamais de compensation : le plancher incompressible).
\end{theoreme}

\begin{proposition}[Pivot barycentre --- \S6.3]\label{prop:bary}
Comme $\mathrm{RR}_p = \mathrm{RR}_0 - m\,x_p$, le pivot $x_p = \mathrm{RR}_0/m =
(\sum_j \nu_j x_j)/m$ --- le barycentre de vega --- annule $\mathrm{RR}_p$ exactement :
le terme de pente disparaît par choix de représentation.
\end{proposition}

\begin{remarque}[Exact contre majorant : autoriser ce que la couche 1 refusait]
Le majorant de couche 1 est $\sum_{j\neq p}\abs{\nu_j}\, d_{j,p}$ (inégalité
triangulaire) : la barre \emph{dans} la somme efface les signes, donc il \textbf{compte
deux fois ce qui s'annule}. Le calcul exact \eqref{eq:tl-exact} peut donc autoriser une
fusion que le majorant refusait --- jamais l'inverse.
\end{remarque}

\paragraph{Décisions, jamais tout-ou-rien (N1--N4).} Les groupes d'une tranche
partagent $\Delta S, \Delta C$ : les coefficients RR/FLF sont \emph{sommés entre groupes
écrasés avant} d'être élevés au carré (la TE globale n'est pas la somme des variances de
groupe). On part tout en \emph{collapse} sur le pivot barycentre ; tant que la TE
globale dépasse le budget, on \emph{extrait} le pire groupe (son $m$ reste netté,
RR/FLY provisionnés en \emph{add-up} avec $s_{\text{skew}}, s_{\text{fly}}$) ; si l'idio
seul déborde, on \emph{scinde}.

\begin{algorithm}[H]
\caption{Couche 2 --- décisions book (N1--N4) sur la structure de couche 1}
\KwEntree{groupes (couche 1), vegas $\nu$, abscisses $x$, modèle de smile, $\alpha$, $\kappa$}
\KwSortie{décisions par groupe, $\mathrm{TE}^2$ exacte, AVA, verdicts}
\KwInit{tous les groupes en \emph{collapse} sur leur pivot barycentre (Prop.~\ref{prop:bary})}
$B \leftarrow (1-\alpha)\Var(\dPi)$\tcp*{$\Var(\dPi)$ de la tranche sous le modèle}
\While{$\mathrm{TE}^2_{\text{globale}} > B$ \textbf{et} des groupes collapsés ont RR/FLY $\neq 0$}{
  extraire le groupe de plus grande $\Var(R)$\tcp*{$m$ netté, RR/FLY en add-up}
}
\While{$\mathrm{TE}^2_{\text{globale}} > B$ \textbf{et} un groupe a $>1$ point}{
  scinder en deux le pire groupe (chaque moitié reprend son pivot barycentre)\;
}
$\mathrm{AVA} \leftarrow \sum_{\text{groupes}} \kappa\big(\abs{m}\,s_{\text{pivot}}
   + [\text{extrait}]\,(\abs{\mathrm{RR}}s_{\text{skew}} + \abs{\mathrm{FLY}}s_{\text{fly}})\big)$\;
\Return{décisions, $\mathrm{TE}^2_{\text{globale}}$, AVA, test/plancher}\;
\end{algorithm}

\begin{impl}
\code{twolayer.py} : \code{SmileModel}, \code{fit\_smile\_model}, \code{model\_distance}
(Th.~\ref{th:tl1}), \code{generated\_correlation} (Prop. 2),
\code{tranche\_dendrogram}/\code{cut\_tranche}, \code{group\_projections},
\code{barycenter\_pivot} (Prop.~\ref{prop:bary}), \code{group\_residual\_variance}
(Th.~\ref{th:tl2}), \code{majorant\_residual\_sd}, \code{book\_variance},
\code{evaluate\_book} (N1--N4). Onglet \emph{Deux couches}.
\end{impl}

% =====================================================================
\section{Audit art.~9(5) et mise en œuvre IPV}\label{sec:audit}
% =====================================================================

\subsection{Dossier de preuve}
Le RTS exige une preuve \emph{documentée} que les expositions nettées se compensent
réellement face à l'incertitude. L'outil assemble, par date de calcul, un dossier JSON
contenant : la partition retenue, les paramètres $(\alpha,\kappa,$ pondération$)$, le
test de variance ($\mathrm{TE}^2$, budget, $R^2$ réalisé, verdict), le plancher, la
décomposition d'AVA, les statistiques par ensemble, les résultats de stress et la trace
gloutonne. Chaque dossier porte une \textbf{empreinte SHA-256} du bundle d'entrée trié :
elle lie le résultat aux données exactes (reproductibilité, détection d'altération).

\begin{impl}
\code{reporting.py} : \code{build\_audit\_pack}, \code{\_bundle\_fingerprint},
\code{audit\_json}. Onglet \emph{Audit \& Export}.
\end{impl}

\subsection{Points de mise en œuvre (\S8 de la note)}
\begin{itemize}[nosep,leftmargin=1.2em]
\item \textbf{$s$} : dispersion inter-contributeurs Totem (écart-type ou intervalles
  interquantiles remis à l'échelle 90~\%) ; peu de contributeurs $\Rightarrow$ gonfler
  $s$, pas le lisser.
\item \textbf{$\rho$} : estimer sur les \emph{variations} des marks de consensus, jamais
  sur les niveaux ; historique court $\Rightarrow$ croiser avec $\Sigma^{\mathrm{daily}}$
  (Prop.~5) et retenir le \emph{moins favorable au netting} ; shrinkage recommandé.
\item \textbf{Asymétrie réglementaire} : le test protège contre la \emph{sous}-estimation
  --- en cas de doute, le stress de corrélation doit être adverse au netting.
\item \textbf{Documentation} : conserver \{partition, $R^2$ réalisé, frontière, stress\}
  par date (preuve art.~9(5)).
\item \textbf{Non-linéarité} : \eqref{eq:dpi} est au premier ordre ; pour des books à
  volga/vanna matériels (cliquets, barrières), valider que les termes de second ordre ne
  réordonnent pas les ensembles --- sinon netter sur scénarios complets plutôt que sur
  sensibilités.
\end{itemize}

% =====================================================================
\appendix
\section{Correspondance théorie $\leftrightarrow$ code $\leftrightarrow$ onglet}
% =====================================================================

\begin{center}\small
\begin{tabular}{@{}p{4.5cm}p{5.2cm}p{4.6cm}@{}}
\toprule
\textbf{Théorie (\S)} & \textbf{Module / fonction} & \textbf{Onglet} \\
\midrule
Modèle, Propriété~\ref{prop:property1}, AVA \eqref{eq:ava-brut}--\eqref{eq:ava-full}
  & \code{model.py}, \code{datasource.py} & Données de marché \\
Passage \eqref{eq:sandwich}, Th.~\ref{th:passage}, \S3.5
  & \code{passage.py} & Passage \\
Netting Déf.~\ref{def:netting}, test \eqref{eq:test}, plancher \eqref{eq:floor}
  & \code{netting.py} & Netting optimal / Labo \\
Th.~\ref{th:two} (deux nœuds)
  & \code{netting.py:two\_bucket} & Labo scénarios \\
Glouton \eqref{eq:opt}, stress \eqref{eq:stress}, robuste
  & \code{optimizer.py} & Netting optimal \\
Distance \eqref{eq:base-risk}, Prop.~\ref{prop:prop8}, Runs 1/2, \S7.4
  & \code{clustering.py} & Structure \& Stabilité \\
Facteurs Prop.~\ref{prop:spectral}, Th.~\ref{th:spectral}, $K^\star$ \eqref{eq:kstar}
  & \code{spectral.py} & Spectral \& Smile \\
Décomposition smile Th.~\ref{th:smile}
  & \code{scenario.py:smile\_decomposition} & Spectral \& Smile \\
Deux couches Déf.~\ref{def:smile}, Th.~\ref{th:tl1}--\ref{th:tl2}, N1--N4
  & \code{twolayer.py} & Deux couches \\
Dossier art.~9(5)
  & \code{reporting.py} & Audit \& Export \\
\bottomrule
\end{tabular}
\end{center}

\section{Notations}
\begin{center}\small
\begin{tabular}{@{}ll@{}}
\toprule
$N,\ \nu$ & matrice $(M,K)$ et vegas par bucket ($\partial V/\partial\sigma$) \\
$\dsig,\ s,\ \rho$ & choc d'incertitude, écarts-types, corrélations \\
$\Sigma = D\rho D$ & covariance des chocs ($D=\operatorname{diag}(s)$) \\
$\kappa,\ \alpha$ & multiplicateur de quantile ($\approx1{,}2816$), seuil du test \\
$\Nt = A_T^\top N A_K$ & vega projeté ; $A_T,A_K$ matrices de passage \\
$\mathcal P,\ m_r,\ W_r,\ \st_r$ & partition, exposition nettée, poids, incertitude repr. \\
$\mathrm{TE}^2,\ R^2$ & tracking error et score de variance \\
$d_{ij},\ \varepsilon$ & risque de base \eqref{eq:base-risk}, hauteur de coupe \\
$\lambda_\ell, u_\ell, c_\ell$ & valeurs/vecteurs propres de $\Sigma$, projections du vega \\
$\mathrm{RR},\ \mathrm{FLY}$ & risk-reversal et butterfly nets (smile) \\
\bottomrule
\end{tabular}
\end{center}

\bigskip\noindent\hrulefill\par
\footnotesize
\noindent\emph{Note de référence du \emph{Vega Netting Studio}. Données de
démonstration synthétiques ; ceci n'est ni un conseil en investissement ni une position
réglementaire. Compilation : \code{pdflatex vega\_netting\_note.tex} (deux passes) ou
Overleaf.}

\end{document}
